\section*{Bloque B. Aplicaciones cuantitativas y financieras}

\noindent\textit{Preguntas 31 a 65. Evalúan la capacidad matemática para calcular tasas efectivas, anualidades, ratios financieros, flujos de caja libres del proyecto (FCLP), VAN, TIR, tasa de Fisher, Costo Anual Equivalente (CAE), WACC corporativo, modelo CAPM y cronogramas de amortización de deuda.}

\q{Un banco comercial aprueba un préstamo a una EPS a una Tasa Nominal Anual del $TNA = 12\%$ con capitalización mensual de intereses. ¿Cuál es la Tasa Efectiva Anual (TEA) equivalente que realmente pagará la empresa?}
\begin{opts}
    \item $TEA = 13.55\%$
    \item $TEA = 11.45\%$
    \item $TEA = 12.00\%$
    \item $TEA \approx 12.68\%$
\end{opts}

\q{Si un depósito a plazo fijo en una entidad financiera rinde una tasa de interés nominal del $i = 8.15\%$ anual y la inflación proyectada para el año es del $\pi = 3.0\%$, la tasa de interés real ($r$) exacta según la ecuación de Fisher es:}
\begin{opts}
    \item $r = 5.15\%$
    \item $r = 11.39\%$
    \item $r = 2.44\%$
    \item $r = 5.00\%$
\end{opts}

\q{Una Planta de Tratamiento genera ahorros netos de energía constantes de 50000 soles anuales que se proyectan a perpetuidad. Si la tasa de descuento corporativa relevante es del $r = 8.0\%$ anual, ¿cuál es el Valor Presente ($VP$) de dichos ahorros?}
\begin{opts}
    \item $VP = 625000$ soles
    \item $VP = 500000$ soles
    \item $VP = 400000$ soles
    \item $VP = 800000$ soles
\end{opts}

\q{Un proyecto hidroeléctrico que abastece a una EPS proyecta un flujo de caja de $C_1 = 30000$ soles para el primer año, el cual crecerá a una tasa anual constante del $g = 3.0\%$ indefinidamente. Si la tasa de descuento es $r = 9.0\%$, el valor presente de la perpetuidad creciente es:}
\begin{opts}
    \item $VP = 1000000$ soles
    \item $VP = 500000$ soles
    \item $VP = 333333$ soles
    \item $VP = 600000$ soles
\end{opts}

\q{Una EPS adquiere un equipo de laboratorio mediante un préstamo de 100000 soles a amortizar en 5 cuotas anuales fijas vencidas iguales a una tasa de interés del $10\%$ anual. El factor de recuperación del capital es $\left[\frac{0.10(1.10)^5}{(1.10)^5 - 1}\right] \approx 0.263797$. La cuota fija anual ($C$) a pagar es:}
\begin{opts}
    \item $C = 24500.50$ soles
    \item $C = 26379.75$ soles
    \item $C = 20000.00$ soles
    \item $C = 30000.00$ soles
\end{opts}

\q{El Balance General de una empresa prestadora registra: Activo Corriente $= 160000$ soles (de los cuales $60000$ soles son Inventarios de repuestos) y Pasivo Corriente $= 80000$ soles. Los ratios de Razón Corriente y Prueba Ácida son respectivamente:}
\begin{opts}
    \item Razón Corriente $= 1.5$ veces; Prueba Ácida $= 1.0$ veces
    \item Razón Corriente $= 2.5$ veces; Prueba Ácida $= 1.75$ veces
    \item Razón Corriente $= 2.0$ veces; Prueba Ácida $= 1.25$ veces
    \item Razón Corriente $= 1.8$ veces; Prueba Ácida $= 0.8$ veces
\end{opts}

\q{Una EPS reporta Activos Totales por 500000 soles y Pasivos Totales exigibles por 300000 soles. ¿Cuál es su Patrimonio Neto y su ratio de Apalancamiento Deuda/Patrimonio ($D/E$)?}
\begin{opts}
    \item Patrimonio $= 200000$ soles y $D/E = 1.50$ (150\%)
    \item Patrimonio $= 150000$ soles y $D/E = 2.00$ (200\%)
    \item Patrimonio $= 200000$ soles y $D/E = 0.60$ (60\%)
    \item Patrimonio $= 300000$ soles y $D/E = 1.00$ (100\%)
\end{opts}

\q{Una empresa reporta Utilidad Operativa antes de intereses e impuestos ($EBIT$) de 120000 soles y registra gastos por intereses de su deuda financiera de 30000 soles en el año. Su ratio de Cobertura de Intereses es:}
\begin{opts}
    \item Cobertura $= 4.00$ veces
    \item Cobertura $= 0.25$ veces
    \item Cobertura $= 5.00$ veces
    \item Cobertura $= 3.00$ veces
\end{opts}

\q{Mediante el sistema DuPont, una empresa presenta: Margen de Utilidad Neta $= 8.0\%$, Rotación de Activos Totales $= 1.5$ veces y Multiplicador del Capital (Activo/Patrimonio) $= 1.50$. ¿Cuál es el $ROE$ de la empresa?}
\begin{opts}
    \item $ROE = 15.0\%$
    \item $ROE = 9.5\%$
    \item $ROE = 18.0\%$
    \item $ROE = 12.0\%$
\end{opts}

\q{Una EPS registra en su gestión operativa: Periodo Promedio de Cobro a usuarios $= 45$ días, Días de Rotación de Inventarios en almacén $= 60$ días y Periodo Promedio de Pago a sus proveedores químicos $= 40$ días. ¿Cuál es su Ciclo de Conversión del Efectivo ($CCE$)?}
\begin{opts}
    \item $CCE = 75$ días
    \item $CCE = 65$ días
    \item $CCE = 55$ días
    \item $CCE = 145$ días
\end{opts}

\q{En la evaluación de un proyecto de microcuencas, la Utilidad Operativa proyectada para el año 1 es $EBIT = 200000$ soles. La tasa de impuesto a la renta es $T = 30\%$, la depreciación del periodo es $40000$ soles, la inversión en activos fijos es $CAPEX = 50000$ soles y la variación en capital de trabajo neto es $\Delta CTN = 10000$ soles. El Flujo de Caja Libre del Proyecto ($FCLP$) para el año 1 es:}
\begin{opts}
    \item $FCLP = 140000$ soles
    \item $FCLP = 180000$ soles
    \item $FCLP = 120000$ soles
    \item $FCLP = 100000$ soles
\end{opts}

\q{Un proyecto de sectorización de redes exige una inversión inicial de $I_0 = 100000$ soles en $t=0$ y genera flujos de caja netos de $40000$ soles al final de cada uno de los siguientes 3 años ($t=1, 2, 3$). La tasa de descuento del costo de capital es $r = 10\%$ anual (factor de anualidad $A_3^{10\%} \approx 2.48685$). El Valor Actual Neto ($VAN$) del proyecto es:}
\begin{opts}
    \item $VAN = +20000$ soles
    \item $VAN = -526$ soles
    \item $VAN = 0$ soles
    \item $VAN = +526$ soles
\end{opts}

\q{Un proyecto pequeño requiere una inversión inicial de $I_0 = 10000$ soles y genera flujos de $6000$ soles en el año 1 y $6000$ soles en el año 2. La Tasa Interna de Retorno ($TIR$) del proyecto satisface $10000 = \frac{6000}{1+TIR} + \frac{6000}{(1+TIR)^2}$. El valor aproximado de la TIR es:}
\begin{opts}
    \item $TIR \approx 20.00\%$
    \item $TIR \approx 13.07\%$
    \item $TIR \approx 10.00\%$
    \item $TIR \approx 16.67\%$
\end{opts}

\q{Considere dos proyectos de inversión mutuamente excluyentes: Proyecto A ($I_0 = -1000$, flujos: $t_1 = 700, t_2 = 700$) y Proyecto B ($I_0 = -1000$, flujos: $t_1 = 200, t_2 = 1300$). La 'Tasa de Fisher' que iguala el VAN de ambos proyectos se obtiene resolviendo $VAN(A) - VAN(B) = 0$. Esta tasa de corte es:}
\begin{opts}
    \item $r_{Fisher} = 10.0\%$
    \item $r_{Fisher} = 8.0\%$
    \item $r_{Fisher} = 20.0\%$
    \item $r_{Fisher} = 15.0\%$
\end{opts}

\q{Un sistema de bombeo solar cuesta 70000 soles y genera ahorros de combustible de 25000 soles anuales constantes durante 5 años. El Periodo de Recuperación simple ($Payback$) del capital es:}
\begin{opts}
    \item $Payback = 4.00$ años
    \item $Payback = 3.50$ años
    \item $Payback = 2.80$ años
    \item $Payback = 2.00$ años
\end{opts}

\q{En la evaluación de un proyecto de control de pérdidas de agua, el Valor Presente de los beneficios económicos totales asciende a $VP(B) = 280000$ soles y el Valor Presente de los costos totales asciende a $VP(C) = 200000$ soles. La Relación Beneficio/Costo ($B/C$) y el $VAN$ son:}
\begin{opts}
    \item $B/C = 1.20$ y $VAN = +40000$ soles
    \item $B/C = 0.71$ y $VAN = -80000$ soles
    \item $B/C = 1.50$ y $VAN = +100000$ soles
    \item $B/C = 1.40$ y $VAN = +80000$ soles
\end{opts}

\q{Para renovar un sistema de telemetría de reservorios se evalúan dos bombas: Bomba Alfa tiene un Valor Presente de Costos de $50000$ soles y vida útil de 3 años ($A_3^{10\%} = 2.48685$). Bomba Beta tiene un Valor Presente de Costos de $75000$ soles y vida útil de 5 años ($A_5^{10\%} = 3.79079$). Con una tasa de descuento del $10\%$, compare sus Costos Anuales Equivalentes ($CAE$):}
\begin{opts}
    \item $CAE_{Alfa} = 15000.0$ soles vs $CAE_{Beta} = 18000.0$ soles
    \item $CAE_{Alfa} = 20105.7$ soles vs $CAE_{Beta} = 19784.8$ soles
    \item $CAE_{Alfa} = 25000.0$ soles vs $CAE_{Beta} = 22000.0$ soles
    \item $CAE_{Alfa} = 21500.0$ soles vs $CAE_{Beta} = 21500.0$ soles
\end{opts}

\q{Una empresa de servicios se financia con $60\%$ de capital propio ($E/V = 0.60$) y $40\%$ de deuda bancaria ($D/V = 0.40$). El costo del capital propio es $K_e = 12.0\%$, la tasa de interés de la deuda es $K_d = 7.0\%$ y la tasa del impuesto a la renta es $T = 30\%$. ¿Cuál es el Costo Promedio Ponderado de Capital ($WACC$)?}
\begin{opts}
    \item $WACC = 9.16\%$
    \item $WACC = 11.20\%$
    \item $WACC = 8.50\%$
    \item $WACC = 10.00\%$
\end{opts}

\q{Aplicando el modelo CAPM adaptado a mercados emergentes ($K_e = R_f + \beta (R_m - R_f) + RP$): la tasa libre de riesgo es $R_f = 4.0\%$, la prima de riesgo de mercado es $(R_m - R_f) = 6.0\%$, la Beta apalancada del sector es $\beta = 0.75$ y el spread de Riesgo País (EMBI+) es $RP = 1.50\%$. El costo del capital propio ($K_e$) es:}
\begin{opts}
    \item $K_e = 11.50\%$
    \item $K_e = 9.25\%$
    \item $K_e = 10.00\%$
    \item $K_e = 8.50\%$
\end{opts}

\q{Una corporación mantiene una deuda financiera perpetua de $D = 3000000$ soles a una tasa de interés del $6.0\%$. Si la tasa del impuesto a la renta corporativa es $T = 30\%$, ¿cuál es el Valor Presente del Escudo Fiscal ($VP_{EF}$) según Modigliani y Miller?}
\begin{opts}
    \item $VP_{EF} = 90000$ soles
    \item $VP_{EF} = 900000$ soles
    \item $VP_{EF} = 180000$ soles
    \item $VP_{EF} = 54000$ soles
\end{opts}

\q{En el sistema de amortización francés de una deuda de 100000 soles a 4 años con cuotas anuales fijas de 31547.08 soles al $10\%$ de interés anual, ¿cuánto corresponde a intereses y cuánto a amortización del principal en la primera cuota?}
\begin{opts}
    \item Interés $= 10000.00$ soles y Amortización $= 21547.08$ soles
    \item Interés $= 3154.71$ soles y Amortización $= 28392.37$ soles
    \item Interés $= 5000.00$ soles y Amortización $= 26547.08$ soles
    \item Interés $= 15000.00$ soles y Amortización $= 16547.08$ soles
\end{opts}

\q{En el sistema de amortización constante alemán de una deuda de 60000 soles a pagar en 3 cuotas anuales al $10\%$ de interés anual sobre saldos:}
\begin{opts}
    \item todas las cuotas totales a pagar resultan idénticas durante los tres años.
    \item no se devengan intereses financieros hasta la cancelación de la última cuota.
    \item la amortización del principal es creciente mientras el interés permanece fijo.
    \item la amortización del principal es fija ($20000$ soles) y las cuotas decrecen.
\end{opts}

\q{Un bono corporativo cupón cero con valor nominal de 1000 soles vence dentro de exactamente 2 años. Si la tasa de rendimiento requerida al vencimiento ($YTM$) es del $8\%$ anual, el precio de mercado actual del bono es:}
\begin{opts}
    \item $P = 920.00$ soles
    \item $P = 857.34$ soles
    \item $P = 840.00$ soles
    \item $P = 1000.00$ soles
\end{opts}

\q{Si un bono tiene una Duración Modificada de $D_{mod} = 4.5$ años y la tasa de interés de mercado sube en 50 puntos básicos ($+0.50\%$), la variación porcentual aproximada en el precio del bono es:}
\begin{opts}
    \item $+0.50\%$
    \item $+2.25\%$
    \item $-2.25\%$
    \item $-4.50\%$
\end{opts}

\q{Un proyecto presenta una inversión de 1000 soles y genera flujos de $600$ soles en año 1 y $800$ soles en año 2. Con una tasa de reinversión del $10\%$, el Valor Futuro terminal en el año 2 es $VF = 600(1.10) + 800 = 660 + 800 = 1460$ soles. La Tasa Interna de Retorno Modificada ($TIRM$) a 2 años es:}
\begin{opts}
    \item $TIRM = 10.00\%$
    \item $TIRM \approx 20.83\%$
    \item $TIRM = 23.00\%$
    \item $TIRM = 15.50\%$
\end{opts}

\q{Una empresa tiene costos fijos operativos de 120000 soles al año. Vende cada metro cúbico tratado a $P = 5.0$ soles y su costo variable unitario es $CV = 2.0$ soles. El punto de equilibrio operativo ($Q_{equilibrio}$) en unidades físicas es:}
\begin{opts}
    \item $Q = 60000 \; m^3$
    \item $Q = 40000 \; m^3$
    \item $Q = 24000 \; m^3$
    \item $Q = 30000 \; m^3$
\end{opts}

\q{El Grado de Apalancamiento Operativo (GAO) de una empresa es $GAO = 2.5$. Si las ventas de la compañía aumentan en un $8.0\%$, ¿en qué porcentaje se incrementará su Utilidad Operativa ($EBIT$)?}
\begin{opts}
    \item $10.5\%$
    \item $16.0\%$
    \item $8.0\%$
    \item $20.0\%$
\end{opts}

\q{Si una firma tiene un Grado de Apalancamiento Operativo de $GAO = 2.0$ y un Grado de Apalancamiento Financiero de $GAF = 1.5$, su Grado de Apalancamiento Total ($GAT$) es:}
\begin{opts}
    \item $GAT = 3.50$
    \item $GAT = 1.33$
    \item $GAT = 0.50$
    \item $GAT = 3.00$
\end{opts}

\q{En el método del Valor Presente Ajustado (APV), un proyecto financiado en parte con deuda subsidiada por el Estado se evalúa sumando:}
\begin{opts}
    \item la TIR ajustada del proyecto multiplicada por el ratio de deuda contable de la firma.
    \item el costo histórico de reposición menos el flujo de caja libre operativo del periodo.
    \item el $VAN_{base}$ sin deuda más el valor presente de escudos fiscales y subsidios.
    \item el beneficio neto contable ajustado por la depreciación y el EBIT sin impuestos.
\end{opts}

\q{Bajo el enfoque de opciones reales, una EPS que diseña una planta de agua con espacio físico disponible para duplicar su capacidad en el futuro si la demanda crece está adquiriendo:}
\begin{opts}
    \item un pasivo contingente de largo plazo que reduce el valor presente neto base.
    \item un contrato forward de cobertura sobre precios de tuberías de PVC y bombeo.
    \item una opción real de abandono financiero equivalente a una opción Put europea.
    \item una Opción Real de Expansión (tipo Call) que añade valor económico al VAN.
\end{opts}

\q{Si el costo del capital desapalancado de una industria es $K_u = 10.0\%$, el costo de la deuda es $K_d = 6.0\%$, la tasa impositiva es $T = 30\%$ y la relación deuda/patrimonio es $D/E = 0.50$, según la Proposición II de Modigliani-Miller con impuestos, el costo de capital propio apalancado $K_e$ es:}
\begin{opts}
    \item $K_e = 12.00\%$
    \item $K_e = 8.60\%$
    \item $K_e = 10.00\%$
    \item $K_e = 11.40\%$
\end{opts}

\q{Al proyectar flujos de caja en presencia de inflación, la regla de consistencia financiera básica exige que:}
\begin{opts}
    \item se excluya por completo la depreciación fiscal de activos de los flujos de caja nominales.
    \item la tasa de descuento nominal equivalga exactamente al doble de la tasa de inflación anual.
    \item los flujos nominales se descuenten a tasa nominal y los reales a tasa real, dando igual VAN.
    \item los flujos de caja nominales se descuenten siempre con la tasa de descuento real sin ajuste.
\end{opts}

\q{Al término del horizonte de evaluación de 10 años de una planta de tratamiento, el Valor Residual (o valor terminal) estimado como perpetuidad del flujo normalizado del año 10 ($FCLP_{10} = 80000$ soles) descontado a una tasa de $r = 10\%$ sin crecimiento ($g=0$) es:}
\begin{opts}
    \item $VR_{10} = 0$ soles al culminar la vida útil en $t=10$
    \item $VR_{10} = 80000$ soles en $t=10$, descontado a $t=0$
    \item $VR_{10} = 400000$ soles en $t=10$, descontado a $t=0$
    \item $VR_{10} = 800000$ soles en $t=10$, descontado a $t=0$
\end{opts}

\q{Si una acción de una empresa de servicios cotiza a 30 soles y reporta una utilidad neta por acción de 2.50 soles, su ratio Precio/Beneficio (P/E) es:}
\begin{opts}
    \item $P/E = 12.0$ veces
    \item $P/E = 8.33$ veces
    \item $P/E = 5.0$ veces
    \item $P/E = 75.0$ veces
\end{opts}

\q{Un bono cero cupón a 3 años de valor nominal 1000 soles se cotiza hoy en el mercado a 751.31 soles. ¿Cuál es su Rendimiento al Vencimiento ($YTM$)?}
\begin{opts}
    \item $YTM = 10.00\%$
    \item $YTM = 12.00\%$
    \item $YTM = 8.00\%$
    \item $YTM = 6.50\%$
\end{opts}

